\section{Thermodynamic formulation and conventions}
\label{sec:model}
% P:model-01
A fluid phase $p$ is described by its species amounts $\mathbf{n}_p$ in mol, its total amount $N_p=\sum_i n_{p,i}$ and its volume $V_p$ in \si{\cubic\metre}.
Its mechanical Helmholtz energy is the ideal density and mixing term plus the residual term~\citep{Gross2001},
\begin{equation}
\frac{A_{p}}{RT}=\sum_i n_{p,i}\left[\ln\frac{n_{p,i}}{V_p\rho_0}-1\right]+N_p\,\ares\!\left(T,\frac{N_p}{V_p},\frac{\mathbf{n}_p}{N_p}\right),
\label{eq:phase-helmholtz}
\end{equation}
where $\rho_0=1$~\si{\mol\per\cubic\metre} is a fixed reference molar density and $\ares$ is the reduced residual molar Helmholtz energy at temperature $T$, molar density $\rho$ and mole fractions $\mathbf{x}$.
Ions are counted as separate species, so $\mathbf{x}$ of an electrolyte solution is a species, not a formula-unit, composition.
Thermodynamic properties and the chemical and mechanical equilibrium equations follow from derivatives of Eq.~\eqref{eq:phase-helmholtz} and the reference terms added in Section~\ref{sec:references}.
Conservation, electroneutrality and normalization supply the remaining equilibrium constraints (Section~\ref{sec:equilibrium}).
Supporting Information Table~\ref{tab:nomenclature} lists the symbols.

% P:model-02
The residual energy is a sum of contributions,
\begin{equation}
\ares=a^{\mathrm{hc}}+a^{\mathrm{disp}}+a^{\mathrm{assoc}}+a^{\mathrm{polar}}+a^{\mathrm{DH}}+a^{\mathrm{Born}},
\label{eq:residual}
\end{equation}
and a parameter record selects which contributions and which permittivity relation enter a calculation (Supporting Information Table~\ref{tab:contributions}).
A contribution that the record does not declare is identically zero, so Eq.~\eqref{eq:residual} is the selected formulation for a calculation, not a fixed inventory.
The hard-chain and dispersion terms are those of Gross and Sadowski, with a binary parameter $k_{ij}$ on the cross dispersion energy that may carry a declared temperature law~\citep{Gross2001,Held2014}.
The optional polar term is either the dipole--dipole contribution of Gross and Vrabec~\citep{Gross2006} or the quadrupole--quadrupole contribution of Gross~\citep{grossEquationofstateContributionPolar2005}; a record declares at most one of them, each is meaningful only with parameters fitted together with it, and no claim is made here for either one combined with ions.

% P:model-03
The association contribution follows Wertheim's theory in the form of Chapman et al.~\citep{Chapman1990},
\begin{equation}
a^{\mathrm{assoc}}=\sum_A x_{i(A)}\nu_A\left(\ln X_A-\frac{X_A}{2}+\frac12\right),\qquad
X_A\Bigl[1+\hat\rho\sum_B x_{i(B)}\nu_B X_B\Delta_{AB}\Bigr]-1=0,
\label{eq:mass-action}
\end{equation}
where $X_A$ is the fraction of sites of type $A$ not bonded and $\nu_A$ the number of such sites on species $i(A)$.
With $\rho$ in \si{\mol\per\cubic\metre}, the number density is $\hat\rho=10^{-30}N_A\rho$ in \AA$^{-3}$ and the association strength $\Delta_{AB}$ is in \AA$^3$, so $\hat\rho\Delta_{AB}$ is dimensionless.
Cross-association parameters follow the combining rule of Gross and Sadowski~\citep{Gross2002} or, when declared, that of Wolbach and Sandler~\citep{wolbachUsingMolecularOrbital1998}.
The site fractions are implicit functions of $(T,\rho,\mathbf{x})$ and of the parameters; Section~\ref{sec:derivatives} differentiates them through the mass-action system.

% P:model-04
The Debye--H\"uckel contribution is~\citep{Held2008,Bulow2019}
\begin{equation}
a^{\mathrm{DH}}=-\frac{\kappa e^2}{12\pi\epsilon_0\epsilon_r k_BT}\sum_{i}x_iz_i^2\chi(\kappa d_i^{\mathrm{DH}}),\qquad
\kappa^2=\frac{\rho N_A e^2\sum_i x_iz_i^2}{\epsilon_0\epsilon_r k_BT},\qquad
\chi(y)=\frac{3}{y^3}\Bigl[\ln(1+y)-y+\frac{y^2}{2}\Bigr],
\label{eq:dh}
\end{equation}
with charge numbers $z_i$, a separately declared closest-approach diameter $d_i^{\mathrm{DH}}$ and the bulk relative permittivity $\epsilon_r$.
The Born contribution is the corrected solvation-shell form of Figiel et al.~\citep{Figiel2025,Figiel2026Correction},
\begin{equation}
\begin{aligned}
a^{\mathrm{Born}}&=-\frac{e^2}{4\pi\epsilon_0k_BT}\sum_{i:z_i\ne0}x_iz_i^2\left[\frac{1-\epsilon_r^{-1}}{D_i}+c_{\mathrm{d}}\left(1-\epsilon_{r,\mathrm{ion}}^{-1}\right)\left(\frac{1}{d_i^{\mathrm{Born}}}-\frac{1}{D_i}\right)\right],\\
D_i&=d_i^{\mathrm{Born}}\left(1+c_{\mathrm{s}}\frac{f_{\mathrm{mix}}-1}{\lvert z_i\rvert}\right)\quad(z_i\ne0),
\end{aligned}
\label{eq:born}
\end{equation}
where $D_i$ is the outer shell diameter of ion $i$, and $f_{\mathrm{mix}}=\sum_jx_jf_j$ averages solvation factors over all neutral and ionic components, with $f_j=1$ for ions.
Here $\epsilon_{r,\mathrm{ion}}$ is the permittivity of the region between the ion and its shell, and $c_{\mathrm{s}}$ and $c_{\mathrm{d}}$ are fixed coefficients that recover the direct Born term when both are zero.
Both electrolyte contributions vanish identically on an ion-free support.
At zero ionic strength $a^{\mathrm{DH}}$ varies as the three-halves power of the ion amounts, which limits the derivatives available there (Section~\ref{sec:availability}).

% P:model-05
The relative permittivity enters Eqs.~\eqref{eq:dh} and~\eqref{eq:born} and must be declared consistently with the parameter record.
Four relations are available:
\begin{equation}
\epsilon_r=
\begin{cases}
\epsilon_{\mathrm{sf}}\big/\bigl(1+\sum_{i\in\mathcal I}\alpha_ix_i\bigr) & \text{ion suppression~\citep{Zuber2014,Figiel2025}},\\[2pt]
\epsilon_{\mathrm{sf}}\sum_{s\in\mathcal S}x_s+\sum_{i\in\mathcal I}x_i\epsilon_{r,i} & \text{component mixing~\citep{AscaniHeld2021,Bulow2019}},\\[2pt]
1+\rho\,x_{\mathrm{solv}}\,d_{V}\bigl(d_{T}/T-1\bigr) & \text{density--temperature~\citep{Yang2024}},\\[2pt]
1+3y\bigl\{1+y+\bigl[\tfrac{17}{16}I-1\bigr]y^2\bigr\} & \text{dipolar perturbation~\citep{Rueben2024}},
\end{cases}
\label{eq:permittivity}
\end{equation}
Here $\mathcal S$ is the set of neutral solvent species, $\mathcal I$ the set of ions, and $x_{\mathrm{solv}}=\sum_{s\in\mathcal S}x_s$ the total solvent mole fraction.
The salt-free solvent mixture has relative permittivity $\epsilon_{\mathrm{sf}}$; $\alpha_i$ is a dimensionless ion suppression coefficient and $\epsilon_{r,i}$ the relative permittivity assigned to ion $i$ in the component-mixing relation.
The coefficients $d_V$ and $d_T$ are averages of the solvent coefficients weighted by $x_s/x_{\mathrm{solv}}$, in \si{\cubic\metre\per\mol} and K, respectively, so $\rho x_{\mathrm{solv}}d_V$ and $d_T/T$ are dimensionless.
In the dipolar-perturbation relation, $y$ is a reduced dipole density and $I$ a dimensionless correlation integral.
The first two relations do not depend on density and describe condensed phases only; the last two tend to $\epsilon_r=1$ as $\rho\to0$.
Pressure, chemical potentials and caloric properties retain $\partial\epsilon_r/\partial\rho$, $\partial\epsilon_r/\partial T$ and $\partial\epsilon_r/\partial\mathbf{x}$ wherever the selected relation depends on them.
The density dependence also decides where a charged trial phase can be evaluated during a phase search (Section~\ref{sec:stability}).

\subsection{Reference states, reactions and units}
\label{sec:references}
% P:model-06
Nonreactive phase equilibrium at fixed temperature needs only Eq.~\eqref{eq:phase-helmholtz}.
Reaction equilibrium also requires the combinations of standard Gibbs energies constrained by the declared reaction targets (see below); a supplied reaction constant does not require a separate ideal-gas thermal record for every species.
Total caloric properties additionally require sufficient species thermal data to determine the reference contributions to enthalpy, entropy and heat capacity.
For a species with a supplied ideal-gas thermal record, the standard Gibbs energy on the density basis of Eq.~\eqref{eq:phase-helmholtz} is
\begin{equation}
\frac{g^{\circ}_{i,\rho}}{RT}=\frac{h^{\circ}_i(T)}{RT}-\frac{s^{\circ}_i(T,P_0)}{R}+\ln\frac{RT\rho_0}{P_0},\qquad
h^{\circ}_i(T)=\Delta_fh^{\circ}_i(T_{\mathrm{ref}})+\int_{T_{\mathrm{ref}}}^{T}C^{\circ}_{p,i}\,\mathrm{d}\tau,
\label{eq:reference}
\end{equation}
with the entropy integrated likewise from $s^{\circ}_i(T_{\mathrm{ref}},P_0)$.
The ideal-gas heat capacities $C^{\circ}_{p,i}$ are given as constant, polynomial, Shomate~\citep{shomateMethodEvaluatingCorrelating1954} in the convention of the NIST Chemistry WebBook~\citep{linstromNISTChemistryWebBook1997}, NASA seven- and nine-coefficient~\citep{gordonComputerProgramCalculation1994} or DIPPR~107~\citep{alySelfconsistentEquationsCalculating1981} correlations over declared temperature intervals.
Reaction targets and supplied thermal records can together determine reference contributions for species without their own records, provided no undetermined neutral combination remains.
Insufficient thermal data make the affected total caloric quantities unavailable; residual properties and reaction-equilibrium calculations remain available when their own input requirements are met.

% P:model-07
A reaction declares its stoichiometry and exactly one target: a constant affinity, a constant equilibrium constant, or the correlation
\begin{equation}
\ln K(T)=a+\frac{b}{T}+c\,\ln\frac{T}{T_{\mathrm{ref}}}+d\,T
\label{eq:lnK}
\end{equation}
on the density basis, for which the activity of species $i$ is $(\rho x_i/\rho_0)\exp(\mu_i^{\mathrm{res}}/RT)$.
A constant on an ideal-gas pressure basis, or one shifted by a standard state with constant reaction enthalpy and entropy, converts exactly into another correlation of the form of Eq.~\eqref{eq:lnK}.
A constant on a molality, infinite-dilution or pure-liquid basis is converted through the residual chemical potentials of the reference state, computed for one pure neutral solvent and a charge-neutral dilution path at each temperature and pressure of the solve; this conversion is called reference completion below.
Only charge-neutral combinations of ionic reference potentials are determined, so single-ion standard states remain a gauge choice.
A reaction constant constrains the standard Gibbs energies only along the reaction subspace; it does not by itself fix the absolute caloric properties of each species, which also require thermal anchors from Eq.~\eqref{eq:reference}.
Thermochemistry that is supplied together with a reaction target is checked against it in Gibbs energy, enthalpy and heat capacity, and an inconsistency is reported as an error rather than resolved in favor of either input.

% P:model-08
Parameters are explicit physical inputs in SI molar units: K, Pa, mol, \si{\cubic\metre} and \si{\joule\per\mol}, with molality in mol per kg of solvent.
Before a calculation the package admits only inputs in the mathematical domain of the selected equations: positive diameters, a declared site topology, records that the selected contributions consume, and a feed whose mole fractions satisfy $\lvert\sum_iz_ix_i\rvert\le10^{-12}$.
A record that no selected equation consumes is refused rather than ignored.
Admission says nothing about whether published parameters are valid at the requested state; that physical validity is the caller's declaration.

% P:model-10
Electrolyte coefficients use the declared molality standard state.
The neutral mean ionic combination obeys $\ln\gamma_\pm=(\nu_+\ln\gamma_++\nu_-\ln\gamma_-)/(\nu_++\nu_-)$, while solvent activity gives $\Phi=-\ln a_w/(M_w\sum_j\nu_jm_j)$ on the stated solvent-mass basis~\citep{HamerWu1972,Held2014}.
Here $M_w$ is solvent molar mass in kg/mol, $m_j$ formula-unit molality and $\nu_j$ its total ion stoichiometry.
Individual-ion conventions cancel only in the neutral combination; mixed-salt isopiestic coefficients retain their source-specific reference convention.

\subsection{Parameter records and their domains}
% P:model-09
A parameter set is part of the physical model.
Adding a polar or solvation contribution while keeping parameters fitted without it changes the model and does not isolate the physical effect of that contribution.
Similarly, changing the diameter in the Debye--H\"uckel or Born terms changes the formulation even when the printed dispersion parameters are identical.
Supporting Information Table~\ref{tab:records} groups the distributed records by their source and scientific role.
The scientific-input collection contains 38 source families, including observation and thermal records; 26 families contain canonical equation-of-state parameter documents.
Several records represent alternative or fitted declarations for the same system, so the number of records is not a count of independently validated fluids.
Source applicability can be incomplete; the temperature and pressure ranges actually compared are given with the results rather than inferred from successful numerical evaluation.
